\section*{Capítulo \ref{ch:b2:gratings} --- Interferência de múltiplas ondas e redes de difração}

\begin{solution}{exo:b2:gratings:1}
$a = \qty{2}{\micro m}$, $\sin\theta_p = 0.275p$: \qty{16.0}{\degree}, \qty{33.4}{\degree}, \qty{55.6}{\degree};
não há quarta ordem. O visível de primeira ordem vai de \qty{11.5}{\degree} a \qty{20.5}{\degree}:
\qty{9}{\degree}. Sob \qty{30}{\degree}: $\sin\theta = 0.5 + 0.275p$: $p = +1$ (\qty{51}{\degree}) de um
lado, $p = -1$ a $-5$ ($13^\circ$, $-3^\circ$, $-19^\circ$, $-37^\circ$, $-61^\circ$) do
outro.
\end{solution}

\begin{solution}{exo:b2:gratings:2}
$a = \qty{1.67}{\micro m}$. Ordem 1: $\theta = \qty{20.7}{\degree}$, $\Delta\theta = \Delta\lambda/a\cos\theta =
\qty{3.9e-4}{rad}$, \qty{0.19}{mm}; ordem 2: $\theta = \qty{45}{\degree}$, $\qty{1.0e-3}{rad}$,
\qty{0.51}{mm}. Pixels de \qty{50}{\micro m} ou menores.
\end{solution}

\begin{solution}{exo:b2:gratings:3}
$N = 12\,000$: $R = 12\,000$, $24\,000$, $36\,000$; $\Delta\lambda = 0.046$, $0.023$, \qty{0.015}{nm};
o dubleto (\qty{0.6}{nm}) e o par H/D (\qty{0.18}{nm}, que exige $R = 3600$)
são separados com folga.
\end{solution}

\begin{solution}{exo:b2:gratings:4}
CD: violeta \qty{14.5}{\degree}, vermelho \qty{25.9}{\degree}; o vermelho existe até a ordem 2, o violeta até a
ordem 4: dois arco-íris (parcialmente superpostos). DVD: violeta \qty{32.7}{\degree}, vermelho
\qty{71}{\degree}; o vermelho só tem a primeira ordem: um arco-íris, mais largo.
\end{solution}

\begin{solution}{exo:b2:gratings:5}
(a) Zeros em $\varphi = 2\pi m/N$, $m = 1, \dots, N - 1$: $N - 1$ zeros e $N - 2$ máximos
entre eles. (b) Em $N\varphi/2 = 3\pi/2$, $I \approx I_0/\sin^2(3\pi/2N) \approx I_0(2N/3\pi)^2$:
$4/9\pi^2 = 0.045$ de $N^2I_0$. (c) $N = 3$: um único \omterm{thm:b2:gratings:nwaves}{máximo secundário} em
$\varphi = \pi$, de altura $I_0$, um nono do principal $9I_0$; $N = 6$: quatro
\omterm{thm:b2:gratings:nwaves}{máximos secundários} de alguns por cento. (d) O lóbulo central de $\sin^2x/x^2$
contém $90\%$ da energia.
\end{solution}

\begin{solution}{exo:b2:gratings:6}
(a) $p > \lambda_{\min}/(\lambda_{\max} - \lambda_{\min}) = 1.33$: a partir de $p = 2$. (b) $\lambda/p$. (c)
\qty{25}{nm}; um prisma cruzado com a rede empilha as ordens uma acima
da outra. (d) $\Delta\lambda = \lambda^2/2L$, uma fração de nanômetro --- centenas
de vezes menor que o de uma rede.
\end{solution}

\begin{solution}{exo:b2:gratings:7}
(a) $\sin\beta = \lambda/2a = 0.3$, $\beta = \qty{17.5}{\degree}$. (b) A faceta isolada difrata
num lóbulo largo centrado em sua direção especular; as ordens que
caem dentro dele recebem a luz. (c) Ordem 2 ($2 \times \qty{500}{nm} = \qty{1}{\micro m}$ no
mesmo ângulo). (d) $R = 2W\sin\beta/\lambda = 3.6 \times 10^5$.
\end{solution}

\begin{solution}{exo:b2:gratings:8}
(a) $\dd\lambda/\dd x = a\cos\theta/pf = \qty{3.1}{nm/mm}$: \qty{50}{\micro m} valem \qty{0.16}{nm}. (b)
$R = 30\,000$: \qty{0.018}{nm}. (c) A fenda, por um fator nove; iguais em $s = \qty{6}{\micro m}$.
(d) A difração na fenda e um detector faminto de luz.
\end{solution}

\begin{solution}{exo:b2:gratings:9}
(a) $p = 2L/\lambda = 16\,700$; ISL $c/2L = \qty{30}{GHz}$, $\lambda^2/2L = \qty{0.036}{nm}$. (b) $\mathcal F
= \pi\sqrt{0.95}/0.05 = 61$; \qty{0.5}{GHz}, \qty{0.6}{pm}. (c) $p\mathcal F = 10^6$, quatro
vezes o da rede de \qty{10}{cm}. (d) Seu \omterm{prop:b2:gratings:fabryperot}{intervalo espectral livre} é uma fração
de nanômetro: as ordens se superpõem a menos que a luz seja pré-filtrada
para uma delas.
\end{solution}

\begin{solution}{exo:b2:gratings:10}
(a) $|\sin\theta - \sin\theta_0| \le 2$. (b) Os raios extremos diferem em caminho de
$W(\sin\theta - \sin\theta_0)$, no máximo $2W$. (c) $2 \times 0.3/5 \times 10^{-7} = 1.2 \times 10^6$; $c/R
= \qty{250}{m/s}$. (d) O centro de uma linha é localizado com um centésimo de sua
largura, e mil linhas baixam a média de $\sqrt{1000}$: metros por
segundo, com uma estabilidade heroica.
\end{solution}

\begin{solution}{exo:b2:gratings:11}
(a) A onda refletida pelo plano inferior percorre $2d\sin\theta$ a mais.
(b) $\sin\theta = 0.273p$: \qty{15.8}{\degree}, \qty{33.1}{\degree}, \qty{55.0}{\degree}. (c) $\lambda \gg 2d$:
nenhum ângulo satisfaz a equação; o vidro não tem planos regulares, daí
apenas anéis difusos --- ele é amorfo. (d) $\sin\theta \le 1$.
\end{solution}

\begin{solution}{exo:b2:gratings:12}
(a) As vizinhas diferem em fase de $2\pi a\sin\theta/\lambda$ pelo caminho e de $-\psi$
pela alimentação; o \omterm{thm:b2:gratings:nwaves}{máximo principal} está em $\varphi = 0$. (b) Primeiro zero em
$\Delta\varphi = 2\pi/N$: $\Delta\theta = \lambda/Na\cos\theta$. (c) Um segundo \omterm{thm:b2:gratings:nwaves}{máximo principal} em
$\varphi = \pm2\pi$ exige $|\sin\theta \pm \lambda/a| \le 1$: isso é impossível para todo ângulo de
apontamento apenas se $a \le \lambda/2$. (d) $Na = \qty{0.96}{m}$: \qty{1.8}{\degree}; $50$ passos,
\qty{50}{\micro s}.
\end{solution}

\begin{solution}{pb:b2:gratings:1}
\textbf{1.} $a = \qty{0.833}{\micro m}$; $N = 144\,000$.

\textbf{2.} $\sin\theta = \sin20^\circ - \lambda/a$ (a ordem fica do outro lado da
normal): \qty{550}{nm} a \qty{18.5}{\degree}, \qty{400}{nm} a \qty{7.9}{\degree}, \qty{700}{nm} a
\qty{29.9}{\degree}: \qty{22}{\degree} de espectro.

\textbf{3.} $1/a\cos\theta = \qty{1.27e-3}{rad/nm}$; $f_2\times$: \qty{1.27}{mm/nm}, ou seja,
\qty{0.79}{nm/mm}, \qty{0.012}{nm} por pixel.

\textbf{4.} $R = N = 144\,000$: \qty{0.0038}{nm}, um terço de pixel.

\textbf{5.} O \qty{400}{nm} de segunda ordem cai onde cairia o \qty{800}{nm} de
primeira ordem, além do vermelho: nenhuma superposição no visível; a superposição
começaria em \qty{350}{nm}, que um filtro de vidro elimina.

\textbf{6.} \qty{22}{\degree} $\times$ \qty{1}{m} $= \qty{38}{cm}$: um detector cobre uma fatia;
gira-se a rede para escolhê-la.

\textbf{7.} $\lambda = 2a\sin17^\circ = \qty{490}{nm}$.

\textbf{8.} $0.8 \times 0.4 = 32\%$.

\textbf{9.} Ângulos \qty{26.5}{\degree}, \qty{14.0}{\degree}, \qty{10.3}{\degree}: $+14$, $-8$ e
\qty{-14}{cm} em relação à posição de \qty{550}{nm}.

\textbf{10.} \qty{0.18}{nm} valem $15$ pixels: separadas.

\textbf{11.} $50$ pixels de separação, cada uma com $4$ pixels de largura.

\textbf{12.} \qty{100}{\micro m}, $6.7$ pixels, \qty{0.08}{nm} --- vinte vezes o
limite da rede: limitado pela fenda.

\textbf{13.} A imagem tem \qty{97}{\micro m} de largura: uma fenda de \qty{50}{\micro m} deixa passar
metade da luz.

\textbf{14.} Menos luz para a mesma nitidez: o ruído de fótons cresce.

\textbf{15.} $\Delta\lambda = \lambda v/c = \qty{0.066}{nm}$, $5.5$ pixels.

\textbf{16.} A velocidade da Terra desloca cada linha de até $\pm5.5$
pixels ao longo do ano; ela é conhecida pelas efemérides muito
melhor que os $0.3\%$ necessários para \qty{100}{m/s}.

\textbf{17.} \qty{2.2e-5}{nm}, $0.002$ pixel --- muito abaixo da largura da linha. Um
centroide com um centésimo de uma linha de \qty{7}{px} dá \qty{0.07}{px}; mil
linhas dão $\sqrt{1000}$: $0.002$ pixel --- o sinal, no limite. Os instrumentos
reais acrescentam poder de resolução e estabilidade.

\textbf{18.} $\Delta\sin\theta = (\lambda/a) \times 10^{-5} = 6.6 \times 10^{-6}$, $\Delta\theta = \qty{7e-6}{rad}$,
\qty{7}{\micro m}: meio pixel por kelvin --- milikelvins para \qty{10}{m/s}.

\textbf{19.} Suas linhas, registradas no mesmo instante sobre os mesmos
pixels, dão a escala de comprimentos de onda e acompanham toda deriva.

\textbf{20.} $a = \qty{3.33}{\micro m}$, $N = 36\,000$, mesmo ângulo (\qty{18.5}{\degree}), $R = 4N
= 144\,000$, mesma dispersão, ISL $\lambda/4 = \qty{137}{nm}$: o poder de resolução e a
dispersão dependem só de $W$ e do ângulo; o \omterm{prop:b2:gratings:fabryperot}{intervalo espectral livre}
encolheu.

\textbf{21.} ISL $\lambda^2/2L = \qty{0.15}{nm}$ ($12$ pixels); $\mathcal F = 30$; picos de
\qty{0.005}{nm} de largura: um pente de linhas finas ao longo do detector, uma
régua de comprimentos de onda.

\textbf{22.} $2 \times 10^{-4}$ contra \qty{1.3e-3}{rad/nm}: seis vezes menos, e o
poder de resolução de um prisma ($b\,\dd n/\dd\lambda \sim 10^4$) e sua escala não linear
também perdem; as redes ainda funcionam no ultravioleta.

\textbf{23.} $R \le 2W/\lambda$: $4 \times 10^6$ para um metro em \qty{500}{nm} --- a largura e
o ângulo são os únicos caminhos para cima.

\textbf{24.} $10^4 \times 0.012 \times 0.32 = 38$ fótons por pixel por segundo; $3800$
em \qty{100}{s}: relação sinal-ruído $\approx 60$.

\textbf{25.} A dispersão, fixada por $p/a\cos\theta$ e pela distância focal da câmera;
o poder de resolução, pela largura iluminada e pelos ângulos ($pN$);
o \omterm{prop:b2:gratings:fabryperot}{intervalo espectral livre}, pela ordem ($\lambda/p$).
\end{solution}
